\chapter{Conclusions and Future Work}\label{chap:conclusions}

This PhD~Thesis, \emph{RicoSLAM: localization and mapping in dynamic
environments}, researched how to perform \gls{slam} with a ground mobile
robot operating in dynamic indoor environments, using wheeled odometry and
2D~laser scanner data as the only sensing modalities. The research was
organized around the development of a corner-like feature detector
(\chapref{chap:features}), a 2D~point cloud registration framework
based on distance maps (\chapref{chap:distance-maps}), and the
RicoSLAM \gls{slam} framework itself (\chapref{chap:slam}), with the
two research questions and the corresponding hypotheses formulated in
\secref{sec:intro:rq+h} as the underlying motivation.
This chapter consolidates the main conclusions drawn from
these contributions and identifies the most relevant directions
for future research.

This chapter is divided into two sections.
First, \secref{sec:conclusions:contributions} summarizes the
conclusions and contributions of the PhD~Thesis.
Then, \secref{sec:conclusions:final-remarks} presents the directions of future
work and final remarks, including the involvement of
the candidate in the scientific community throughout the PhD~Thesis.

\section{Conclusions and Contributions}\label{sec:conclusions:contributions}

The research questions and hypotheses formulated in
\secref{sec:intro:rq+h} guided the development of the contributions
reported in \chapref{chap:features} to \chapref{chap:slam}.
This section is organized in two parts.
First, \secref{sec:conclusions:rq+h} restates both research questions and
reports, hypothesis by hypothesis, the outcome supported by the
experimental evidence gathered throughout this PhD~Thesis.
Then, \secref{sec:conclusions:summary} consolidates the main conclusions
drawn from the corresponding chapters.
In \secref{sec:conclusions:rq+h}, a hypothesis is reported as
\emph{supported} when the experiments conducted in this PhD~Thesis validate
its claim under the conditions stated in the corresponding chapter, and as
\emph{partially supported} when the proposed mechanism was formulated,
implemented, and shown to operate, but either its isolated contribution was
not quantified through a dedicated ablation study or its validation did not
reach end-to-end operation in real-world conditions.

\subsection{Research questions and hypotheses}\label{sec:conclusions:rq+h}

\subsubsection{Research Question~\ref{question:mapping}}

\emph{How to improve the mapping results when performing \gls{slam} with a
ground mobile robot in a dynamic indoor environment?}

The four hypotheses associated with this research question address the
continuous update of the local environment representation
(Hypothesis~\ref{hypothesis:mapping:ray-tracing}), the bounded growth of the
back-end graph (Hypothesis~\ref{hypothesis:mapping:relocalization}),
the removal of dynamic observations from the final map representation
(Hypothesis~\ref{hypothesis:mapping:clean-map}), and the offline refinement
of the global map (Hypothesis~\ref{hypothesis:mapping:ogba}):
\begin{itemize}
\item \textbf{Hypothesis~\ref{hypothesis:mapping:ray-tracing} --
ray tracing to update local maps online: partially supported.}
The ray-based merger proposed in \secref{sec:slam:dynamics:online} maintains
the local map of the current keyframe by adding the returns falling in front
of the mapped point on the same ray and by replacing the map points that the
latest beam sees through, rebuilding the keyframe's distance map after every
update so that the next alignment registers against the most recent state of
the scene. This mechanism supports the local map operating modes of RicoSLAM,
which attained the best online \gls{ate} results of
\tabref{tab:slam:results:iilabs3d} on Nav~A~Diff (\SI{0.047}{\metre}),
Nav~A~Omni (\SI{0.030}{\metre}), and Loop (\SI{0.033}{\metre}), and it keeps
the registration reference consistent with the latest observations on the
industrial sequences of \secref{sec:slam:results:dynamics}.
Nevertheless, the experimental evaluation does not include an ablation study
in which the online merger is enabled and disabled under otherwise identical
conditions. Therefore, the part of the hypothesis claiming an increase of the
inlier ratio and of the pose estimation stability is supported by the
operation of the framework, but it is not quantified in isolation;
\item \textbf{Hypothesis~\ref{hypothesis:mapping:relocalization} --
relocalization into existing graph representations: supported.}
The relocalization policy proposed in
\secref{sec:slam:online:relocalization} transfers the tracker into an
existing keyframe whenever a short-loop candidate is closer to the robot than
the current keyframe, so that revisits reuse the existing graph nodes instead
of creating new ones. On the four planar sequences of the
IILABS~3D~\cite{ribeiro:access:2025} dataset
(\tabref{tab:slam:results:iilabs3d:graph}), the local map mode with
relocalization ends the run with $17$, $17$, $18$, and $8$ keyframe
variables, between $5.8$ and $18$ times fewer nodes than the most compact
SLAM~Toolbox~\cite{macenski-jambrecic:joss:2021} configuration ($310$, $126$,
$261$, and $46$), while its offline corrected \gls{ate}
(\SI{0.043}{\metre}, \SI{0.031}{\metre}, \SI{0.039}{\metre}, and
\SI{0.019}{\metre}) remains at or below the best
SLAM~Toolbox~\cite{macenski-jambrecic:joss:2021} result on every sequence
(\SI{0.053}{\metre}, \SI{0.061}{\metre}, \SI{0.040}{\metre}, and
\SI{0.042}{\metre}). On the Loop sequence, in which the robot repeats
5~cycles through the same space, the graph drops from $83$ to $19$ nodes in
scan mode and from $89$ to $18$ nodes in local map mode once relocalization
is enabled. These results confirm that the back-end size follows the spatial
extent of the explored environment rather than the duration of the operation
or the path length;
\item \textbf{Hypothesis~\ref{hypothesis:mapping:clean-map} --
offline removal of dynamic points during map generation: supported.}
The offline stage proposed in \secref{sec:slam:dynamics:offline} re-evaluates
the same ray model over all the scans retained by the framework, now placed
in the world frame with the optimized keyframe poses, and combines it with a
Bresenham-like~\cite{bresenham:ibm:1965,zingl:uas:2016} rasterization so that
only the surviving endpoints contribute to the occupancy statistics.
Across the five datasets considered in \secref{sec:slam:results:dynamics}
(IILABS~3D~\cite{ribeiro:access:2025}, Avantgarde, Cappero, Flowbotic, and
Kuka), and for every sequence, operating mode, and threshold evaluated, the
proposed mechanism sharpened the empirical occupancy distribution of the
resulting probability map relative to the \emph{standard} rasterization
baseline, while reclassifying to an \emph{unknown} state the cells whose only
supporting observations were transient dynamic returns.
The single-return ray model bounds this result: a return is only discarded
once a later beam traverses the same ray and sees through it, leaving the
semi-static changes observed on the Avantgarde sequences unresolved;
\item \textbf{Hypothesis~\ref{hypothesis:mapping:ogba} --
offline global map refinement: partially supported, demonstrated in
principle.}
The distance map-based \gls{gba} factors proposed in \secref{sec:slam:gba}
replace the aggregated $SE(2)$ pose-pose measurements by per-point residuals
evaluated against the distance map of a neighboring keyframe, with their
Jacobians derived via right-sided perturbation on the $SE(2)$ manifold.
The synthetic benchmark of \secref{sec:slam:results:gba} showed these factors
recovering a Manhattan-like graph of $289$~poses perturbed by
$\pm\SI{1.5}{\metre}$ on translation and $\pm\SI{20}{\degree}$ on
orientation, both under noise-free and $\sigma = \SI{0.03}{\metre}$ range
noise conditions, with $\mathcal{D}_\mathrm{p2p}(\mathcal{M}_1)$ and
voxelization attaining the lowest errors under noise.
However, that benchmark supplies the factor graph instead of building it:
the graph topology follows the ground-truth proximity between poses, each
variable holds a single clean scan of a static scene, and the point clouds
are already expressed in a consistent reference frame.
Under real operation, the same inputs must come from the online pipeline,
where keyframe local maps only partially overlap, may still contain dynamic
returns, and the neighborhood of each variable must be determined from
estimated rather than ground-truth poses.
Consequently, the error formulations of the offline global refinement are
validated and remain valid for real-world data, since neither the residuals
nor their Jacobians depend on the synthetic nature of the scene, but the
effect of the \gls{gba} stage on the global map produced by RicoSLAM in
real-world deployments is not yet demonstrated
(see \secref{sec:conclusions:future-work}).
\end{itemize}

Summing up, the mapping results of \gls{slam} in dynamic indoor environments
improve by keeping the reference representation consistent with the latest
observations, by reattaching the tracker to the existing graph on revisits
instead of extending it, and by rejecting the dynamic returns through ray
tracing before they reach the final occupancy grid map. The offline global
refinement completes this answer at the formulation level, while its
end-to-end validation remains open.

\subsubsection{Research Question~\ref{question:localization}}

\emph{How to integrate environment dynamics in the localization process when
performing \gls{slam} with a ground mobile robot in a dynamic indoor
environment?}

The three hypotheses associated with this research question address the
extraction of geometric features from raw 2D~laser scanner data
(Hypothesis~\ref{hypothesis:localization:features}), the use of distance maps
as a correspondence-free registration representation combined with robust
estimators (Hypothesis~\ref{hypothesis:localization:distance-maps}), and the
continuous update of local maps for a more stable scan-to-map localization
(Hypothesis~\ref{hypothesis:localization:local-maps}):
\begin{itemize}
\item \textbf{Hypothesis~\ref{hypothesis:localization:features} --
geometric feature extraction from 2D~laser scanner data: supported at the
feature extraction level.}
The corner-like detector proposed in \chapref{chap:features} extracts stable
geometric information from raw 2D~laser data: $6$~features attained a
$100\%$ observability ratio under the sensor \gls{fov} and the pose
uncertainty of the experiments, $20$ and $16$ features remained observable
above $50\%$ and $75\%$, respectively, and the observability ratio elsewhere
in the map stayed below $1.0\%$, indicating few false positives.
Regarding geometric invariance, $91.5\%$ of the possible associations among
the $39$~most persistent features showed a standard deviation at or below
\SI{0.03}{\metre}, matching the noise of the Hokuyo~UST-10LX sensor.
The detector was nonetheless not adopted as the front-end of RicoSLAM, as a
deliberate scoping decision: a detector is not yet an estimator, since
per-feature covariance, feature descriptors, and a data association policy
would still be required (see \secref{sec:conclusions:future-work}); geometric
features do not remove the dynamics problem by themselves; and the distance
map representation of \chapref{chap:distance-maps} already provides data
association and cost evaluation at $O(1)$ over raw scans.
Thus, the hypothesis is supported as far as extracting stable and viewpoint
invariant geometric information from 2D~lasers is concerned, with the
integration into a feature-based tracker left as future work;
\item \textbf{Hypothesis~\ref{hypothesis:localization:distance-maps} --
distance maps for 2D~point cloud registration: supported.}
The framework of \chapref{chap:distance-maps} replaces the nearest-neighbor
search performed at every \gls{icp}~\cite{besl-mckay:tpami:1992,chen-medioni:icra:1991}
iteration by up to $O(1)$ grid lookups, supporting point-to-point and
point-to-plane residuals, three representation levels, dense or sparse
storage, and the classical \gls{icp} baselines under a single data structure.
The precomputed point-to-point variant
$\mathcal{D}_\mathrm{p2p}(\mathcal{M}_1)$ lowered the odometric drift
relative to both of its analytical point-to-point counterparts in most
configurations, reaching drift levels competitive with the point-to-plane
formulations without ever estimating surface normals, and produced less
overconfident information matrices in degenerate scenes.
Dynamic and spurious observations are handled inside the optimization by the
Cauchy robustifier and by the truncation of the distance field at
$d_{\max}$, which bounds the residual and cancels its gradient beyond the
truncation band;
\item \textbf{Hypothesis~\ref{hypothesis:localization:local-maps} --
continuously updated local maps: supported.}
Aggregating the scans into the local map of the current keyframe improved the
online \gls{ate} of RicoSLAM over the corresponding scan-based variants on
Nav~A~Diff (from \SI{0.063}{\metre} to \SI{0.047}{\metre}) and on Nav~A~Omni
(from \SI{0.053}{\metre} to \SI{0.043}{\metre}), while matching it on Loop
(\SI{0.037}{\metre}), as reported in \tabref{tab:slam:results:iilabs3d}.
Beyond trajectory accuracy, the local map is also the representation against
which the loop closure and the relocalization candidates are validated
through a scan-to-map alignment, and the one on which the online dynamics
removal operates. The exception is the Slippage sequence
(\SI{0.015}{\metre} in scan mode against \SI{0.019}{\metre} in local map
mode), where the short straight-line trajectory gives the merger little
opportunity to improve the reference while still accumulating the
registration residuals of the merged scans.
\end{itemize}

Summing up, environment dynamics are integrated in the localization process
by registering every scan against a reference that is continuously updated
with the latest observations, and by treating the remaining dynamic returns
as outliers inside the registration itself, through the robust estimators and
the truncation of the distance field.

\subsection{Summary of the contributions}\label{sec:conclusions:summary}

The remainder of this section summarizes the main conclusions drawn from
the corresponding chapters.

\chapref{chap:features} introduces a line fitting-based corner-like
detector for 2D~laser scanner data~\cite{sousa:icara:2024}, in which
the keypoint is computed as the closed-form intersection of two
locally fitted lines at multiple distance-based neighborhood scales. This
formulation avoids the dependency on global line-segment extractors
that characterizes most literature corner detectors and produces
feature locations that do not need to coincide with a particular
laser return. Experimental results showed that the proposed
corner-like features exhibit viewpoint invariance, where
$6$~features attained a $100\%$ observability ratio under the
considered sensor \gls{fov} and pose uncertainty, while $20$ and $16$ features
remained observable above $50\%$ and $75\%$ across the experiments, respectively.
In addition, the analysis of the association distance among the $39$
most persistent features showed that $91.5\%$ of the possible
feature associations have a standard deviation at or below
$\SI{0.03}{\metre}$, consistent with the characteristic noise of the
Hokuyo~UST-10LX 2D~laser scanner sensor while yielding geometric invariance.
These results indicate that the proposed corner-like features are stable enough
to be used as input to a feature-based tracking module for mobile robots.

\chapref{chap:distance-maps} introduces a 2D~point cloud registration
framework that replaces the explicit nearest-neighbor search performed
at every iteration of classical
\gls{icp}~\cite{besl-mckay:tpami:1992,chen-medioni:icra:1991} with up
to $O(1)$ grid lookups on a precomputed distance map data structure.
Both point-to-point and point-to-plane error formulations are derived
within the Gauss-Newton framework on the $SE(2)$ manifold, with the
Jacobians of both variants obtained through right-sided perturbation.
The unified distance map data structure introduced in this chapter
supports dense and sparse storage layouts and three representation
levels~$\mathcal{M}_l$, while remaining compatible with the classical
\gls{icp} formulations used as baselines. Across the synthetic and
IILABS~3D~\cite{ribeiro:access:2025} benchmarks, the precomputed
point-to-point variant $\mathcal{D}_\mathrm{p2p}(\mathcal{M}_1)$ stood
out as the best-performing point-to-point formulation, consistently
improving over its analytical counterparts
$\mathcal{D}_\mathrm{p2p}(\mathcal{M}_0)$ and
$\mathrm{ICP}_{\mathcal{D},\,\mathrm{p2p}}$ on the \gls{rte} and
\gls{rre} metrics across all four planar sequences of the
IILABS~3D~\cite{ribeiro:access:2025} dataset (Nav~A~Diff, Nav~A~Omni,
Loop, and Slippage), in agreement with the synthetic results
discussed in \chapref{chap:distance-maps}. More importantly,
$\mathcal{D}_\mathrm{p2p}(\mathcal{M}_1)$ also reached competitive
accuracy against the analytical point-to-plane formulations
$\mathcal{D}_\mathrm{p2pl}(\mathcal{M}_0)$ and
$\mathrm{ICP}_{\mathcal{D},\,\mathrm{p2pl}}$. On Nav~A~Omni
at $r_\mathrm{max} = \SI{30}{\metre}$, $\mathcal{D}_\mathrm{p2p}(\mathcal{M}_1)$
obtained the lowest drift overall both on \gls{rte}
(\SI{0.81}{\percent}) and on \gls{rre} (\SI{0.063}{\degree\per\metre}).
On Nav~A~Diff at the same range, $\mathcal{D}_\mathrm{p2p}(\mathcal{M}_1)$
attained the lowest \gls{rre} overall (\SI{0.078}{\degree\per\metre}) with
a \gls{rte} of \SI{0.95}{\percent}, within \SI{0.07}{\percent} of the
$\mathrm{ICP}_{\mathcal{D},\,\mathrm{p2pl}}$ baseline. On the Loop sequence,
the point-to-plane formulations retained an advantage, with
$\mathrm{ICP}_{\mathcal{D},\,\mathrm{p2pl}}$ attaining
\SI{1.37}{\percent} \gls{rte} and \SI{0.106}{\degree\per\metre}
\gls{rre} against \SI{2.14}{\percent} and \SI{0.123}{\degree\per\metre}
for $\mathcal{D}_\mathrm{p2p}(\mathcal{M}_1)$, consistent with the
chapter's discussion of the residual gap that emerges when traversing
large spaces through narrow passages. Overall, these results
underline the relevance of the precomputed distance field to
regularize point-to-point correspondences against sensor noise and to
provide an effective alternative to analytical point-to-plane
formulations in a broad range of operating conditions,
avoiding kd-tree search~\cite{bentley:cacm:1975} traditionally required at every
\gls{icp}~\cite{besl-mckay:tpami:1992,chen-medioni:icra:1991} iteration.

\chapref{chap:slam} presents the RicoSLAM framework developed within the scope
of this PhD~Thesis, including the distance map-based tracker from
\chapref{chap:distance-maps} with split and merge triggering criteria for
keyframe management, a loop closure detection procedure based on world-frame
proximity and graph-distance, and a relocalization policy that
transfers the tracker into an existing keyframe of the back-end graph
whenever a short-loop candidate is available. On the
IILABS~3D~\cite{ribeiro:access:2025} dataset, RicoSLAM attains the
lowest \gls{ate} across its operating modes on all four planar
sequences considered in the trajectory accuracy benchmark, with
\SI{0.047}{\metre} on Nav~A~Diff (\emph{local} mode),
\SI{0.030}{\metre} on Nav~A~Omni (\emph{local~w/upd.}\ mode),
\SI{0.033}{\metre} on Loop (\emph{local~w/upd.}\ mode), and
\SI{0.014}{\metre} on Slippage (\emph{scan~w/reloc.}\ mode). These
values are lower than the \SI{0.053}{\metre}, \SI{0.061}{\metre},
\SI{0.040}{\metre}, and \SI{0.042}{\metre} obtained by the best
SLAM~Toolbox~\cite{macenski-jambrecic:joss:2021} configurations, and
substantially lower than the GMapping~\cite{grisetti:tro:2007}
results, which fail to close the loop on the Loop sequence.
The offline pose-correction stream proposed in
\secref{sec:slam:online:pose-history} further reduces the \gls{ate} on
the three longer sequences, while its dispersion analysis suggests
that the residual error is dominated by systematic sources
(\eg, extrinsic calibration, sensor noise, ground-truth precision)
rather than by front-end drift or loop closure misalignment. In
terms of scalability, all RicoSLAM variants produce smaller back-end graphs than
SLAM~Toolbox~\cite{macenski-jambrecic:joss:2021} for comparable, and
often better, trajectory accuracy. The combination of the
relocalization policy with the local-map operating mode
(\emph{local~w/reloc.}\ variant) produces a back-end whose size
depends primarily on the spatial extent of the explored environment
rather than on the duration of operation or the path length, even
when revisiting the same locations multiple times.

Furthermore, the dynamics removal strategy proposed in \chapref{chap:slam}
operates at two complementary stages of RicoSLAM. The online stage
runs during local map maintenance and treats each laser beam as an
observation from the surrounding scene that may either see through a
previously mapped point or add new observations to the local
representation. The offline stage runs during 2D~occupancy grid map
generation and combines the same ray model with a
Bresenham-like~\cite{bresenham:ibm:1965,zingl:uas:2016} line
rasterization to suppress dynamic returns before integrating the
scans into the final occupancy grid. The experimental evaluation
considered the IILABS~3D~\cite{ribeiro:access:2025} dataset together
with four additional datasets from real-world \gls{amr} deployments
(Avantgarde, Cappero, Flowbotic, and Kuka), spanning sensor sampling
rates from \SI{6}{\hertz} to \SI{40}{\hertz}, sequence durations
from a few minutes up to roughly \SI{78}{\minute}, and path lengths
from tens of meters up to kilometers. Across every sequence,
operating mode, and threshold considered in the experiment, the
proposed offline dynamics removal mechanism consistently sharpened
the empirical occupancy distribution of the final probability map
relative to the standard rasterization baseline and reclassified to
an \emph{unknown} state cells whose only supporting observations were
transient dynamic returns. Together with the trajectory accuracy
and scalability results, these mapping experiments demonstrate that
RicoSLAM is able to produce a cleaner and more compact representation
of the environment suitable for downstream localization in
autonomous operation, even in the presence of dynamic returns and
sensor noise from industrial \gls{amr} deployments.

Lastly, the distance map-based \gls{gba} factors introduced in
\secref{sec:slam:gba} augment the pose-pose constraints of the online
graph with per-point distance map factors, both in their
point-to-point and point-to-plane formulations. The Jacobians of
these factors are derived via right-sided perturbation on the
$SE(2)$ manifold, building on the corresponding error formulations
from \chapref{chap:distance-maps}. The synthetic benchmark evaluated
four \gls{gba} factor configurations on a Manhattan-like grid of
$289$ poses perturbed by $\pm\SI{1.5}{\metre}$ on translation and
$\pm\SI{20}{\degree}$ on orientation, considering noise-free
conditions and a Gaussian range noise of
$\sigma = \SI{0.03}{\metre}$. Although the graph topology and the
synthetic scene are regular and well-structured, the results
indicate that the proposed factors are suited to globally refining
the back-end pose graph produced by the online RicoSLAM pipeline.
It is worth emphasizing what this benchmark does and does not assume.
The factor graph is given rather than built. Its topology follows the
ground-truth proximity between poses, every variable holds a single
scan of a static scene, and no association errors can occur between a point
and the distance map of a neighboring variable.
Conversely, nothing in the derivation of the residuals or of their Jacobians
depends on the synthetic nature of the data. So, the proposed formulations
remain valid for real-world observations.
Therefore, the open problem is not the formulation of the factors, but the
construction of the graph that the \gls{gba} stage consumes, with the
partially overlapping local maps, the residual dynamic returns, and the
pose estimates produced by the online pipeline.
Consistent with the registration evaluation of
\chapref{chap:distance-maps}, the precomputed point-to-point variant
$\mathcal{D}_\mathrm{p2p}(\mathcal{M}_1)$ with voxelization of the
moving point clouds attained the lowest translational and rotational
errors under range noise conditions, reinforcing the relevance of
distance map-based point-to-point formulations even when compared
against point-to-plane variants. The full integration of the
proposed \gls{gba} stage with the online RicoSLAM pipeline and its
end-to-end validation in real-world environments remains a work in
progress at the time of writing this PhD~Thesis, and is identified as
the main open direction in \secref{sec:conclusions:final-remarks}.

Beyond the scientific contributions discussed above,
the PhD~Thesis materializes into a set of peer-reviewed publications
covering the systematic literature review on long-term localization
and mapping~\cite{sousa:jfr:2023} (\appxref{appx:slr}), the corner-like feature
detector~\cite{sousa:icara:2024} (\appxref{appx:features}), the multimodal
mobile platform~\cite{sousa:icarsc:2025} (\appxref{appx:hangfaq2})
used for data acquisition of the IILABS~3D~\cite{ribeiro:access:2025}
(\appxref{appx:iilabs3d}), the distance map-based 2D~point cloud registration
framework of \chapref{chap:distance-maps}~\cite{sousa:robot:2026}
(\appxref{appx:dmaps}), and the usage of RicoSLAM during container unloading
applications~\cite{lopes:icara:2026} (\appxref{appx:cargo}).
The RicoSLAM framework of \chapref{chap:slam} is additionally reported in a
manuscript in preparation~\cite{sousa:ral:2026} (\appxref{appx:ricoslam}),
which extends the evaluation of this PhD~Thesis with the
Cartographer~\cite{hess:icra:2016} baseline and with the total execution time
of every system considered.
These publications are complemented by
the public release of the IILABS~3D~\cite{ribeiro:access:2025}
dataset and of the corresponding software repositories listed in
\secref{sec:intro:contributions}. The combination of these
contributions provides a unified \gls{slam} solution that ties
together distance map-based registration,
dynamics-aware local mapping, and offline global refinement,
addressing in an integrated manner the two research questions on
improving the localization and mapping processes in dynamic environments
formulated in \secref{sec:intro:rq+h}.

\section{Future Work and Final Remarks}\label{sec:conclusions:final-remarks}

This PhD~Thesis advanced the state of the art in 2D~\gls{slam} for
dynamic indoor environments along several complementary directions,
from feature extraction and 2D~point cloud registration to dynamics-aware
\gls{slam} and offline global refinement. The main results discussed in
\secref{sec:conclusions:contributions} validate the hypotheses formulated in
\secref{sec:intro:rq+h} and provide a basis for further developments.
Nevertheless, the research also exposes a set of open problems and natural
extensions that were identified during the development of this PhD~Thesis,
but whose execution was beyond its scope. This section consolidates the
directions of future work and provides an overview of the candidate's
involvement in the scientific community throughout the PhD~Thesis.

\subsection{Future work}\label{sec:conclusions:future-work}

The directions of future work identified in
\chapref{chap:features}, \chapref{chap:distance-maps}, and
\chapref{chap:slam} are consolidated as follows, ranging from the front-end
perception modules to the offline global refinement stages:
\begin{itemize}
\item \textbf{Covariance estimation for the corner-like features.}
The corner-like feature detector proposed in \chapref{chap:features}
produces the keypoint location as the closed-form intersection of two
locally fitted lines at multiple distance-based scales, but does not
yet provide an explicit uncertainty model for the feature observation.
Propagating the covariance of the line-fitting parameters through the
closed-form intersection step would yield a per-feature covariance
suitable for weighting the feature observations within a probabilistic
data association and tracking module. This extension is a prerequisite
for the integration of the corner-like features as observations within
a feature-based tracking module;
\item \textbf{Hybrid raw-scan, feature, and beacon-based \gls{slam}
framework.}
The current RicoSLAM framework focuses on the raw 2D~laser scanner data
as the only perceptual modality. A natural extension is the
development of a hybrid pose-landmark \gls{slam} framework that
simultaneously supports raw scan observations, geometric features
(\eg, line segments, the corner-like features proposed in
\chapref{chap:features}), and artificial beacons (\eg, reflective
cylinders). Such a framework would combine the robustness of dense
distance map-based registration with the long-range associability of
beacon-based and feature-based observations, while reusing the
$SE(2)$ graph back-end already adopted by RicoSLAM. The
covariance estimation extension for the corner-like features mentioned
above is a prerequisite for this development;
\item \textbf{Improvements on the precomputed point-to-plane
formulation.}
The precomputed point-to-plane variant
$\mathcal{D}_\mathrm{p2pl}(\mathcal{M}_2)$ achieves competitive
registration accuracy with respect to the \gls{icp} baselines, but
is more sensitive to range noise than the precomputed point-to-point
variant $\mathcal{D}_\mathrm{p2p}(\mathcal{M}_1)$, attributed to the
compounding of finite-difference errors when approximating both the
gradient and the Hessian of the distance field. Alternative strategies
to estimate these quantities, such as semi-analytical formulations that combine
the analytical gradient at $\mathcal{M}_1$ with a smoothed Hessian, or
truncated-distance reformulations that explicitly handle the
truncation boundaries, may reduce this sensitivity and improve the
robustness of the precomputed point-to-plane variant under noise.
\item \textbf{Improvements on the uncertainty estimation analysis for
2D~point cloud registration.}
The pose-error consistency metric
$\chi^2_{\boldsymbol{\varepsilon}}$ adopted in
\secref{sec:distance-maps:results:synthetic} is an
approximation rather than a true Mahalanobis statistic, because the
target-frame pose error
$\boldsymbol{\varepsilon}\!\left(\mathbf{x}\right)$ derived from
$\mathbf{x}\cdot\mathbf{x}^{*,-1}$ and the body-frame information
matrix $\boldsymbol{\Omega}$ returned by the Gauss-Newton solver
under the right-sided perturbation update live on different tangent
spaces of $SE(2)$. Both representations only coincide, up to the
adjoint twist $\mathrm{Ad}_{\mathbf{x}^*}^T\,\boldsymbol{\Omega}\,
\mathrm{Ad}_{\mathbf{x}^*}$, in the small angular residual regime.
Nevertheless, the target-frame parametrization was preferred in this
PhD~Thesis because of its diagnostic value in degenerate scenes,
projecting otherwise unobservable residuals onto the observable directions of
$\boldsymbol{\Omega}$ and inflating the metric whenever the alignment
fails along a degenerate axis.
\item \textbf{Improvements on the online dynamics removal mechanism.}
The online dynamics removal stage of RicoSLAM operates on a
per-keyframe basis and commits the updated local map directly to the
back-end graph once a split or merge criterion is triggered. An
intermediate buffering stage between the tracker and the back-end,
which could accumulate evidence over a sliding window of keyframes
before committing them, may further reduce the influence of transient
returns on the back-end graph and improve the robustness of the
framework in highly dynamic scenarios. In addition, the simplified
single-return ray model adopted by the online and offline dynamics
removal mechanisms could be extended to multi-return formulations,
addressing the residual semi-static dynamics observed in the Avantgarde
sequences and improving the cleaning effect on the final occupancy
grid representation.
The evaluation of this mechanism also deserves a stronger validation
protocol. The occupancy distribution analysis of
\secref{sec:slam:results:dynamics} measures how decided the resulting map
becomes, but a sharper map does not prove that every discarded return was
truly dynamic, nor that no sparsely observed static structure was removed.
A manually labelled or independently surveyed static reference map would
enable reporting precision and recall of the removal decision, the rate of
falsely removed static returns, and the rate of missed dynamic returns,
turning the current evidence into a quantitative assessment of the
mechanism;
\item \textbf{Integration and real-world validation of the distance
map-based \gls{gba}.}
The distance map-based \gls{gba} factors proposed in
\secref{sec:slam:gba} were validated on a synthetic benchmark with
ground-truth poses, isolating their contribution from the online
operation of RicoSLAM. The full integration of the proposed \gls{gba}
stage with the online RicoSLAM pipeline and its end-to-end validation
in real-world environments stands out as the most immediate direction
of future work. This integration would close the loop between the
online graph \gls{slam} operation and the offline global refinement
stage, enabling the explicit measurement of the contribution of the
\gls{gba} stage to the trajectory accuracy and to the quality of the
2D~occupancy grid maps produced by RicoSLAM in real-world industrial
deployments.
The main obstacle for this integration is not the error formulation, which
is independent of the origin of the data, but the construction of the graph
that the \gls{gba} stage optimizes. The synthetic benchmark of
\secref{sec:slam:results:gba} assumes that this graph is already available.
The neighborhood of each variable is derived from the ground-truth poses, and
each variable holds one clean scan of a static scene. So, no association
error can occur between a point and the distance map of a neighboring
variable. Building the same structure from a real run requires selecting
which pairs of keyframes share enough observed surface to be linked, handling
local maps that only partially overlap and that may still carry dynamic or
semi-static returns, initializing the optimization from estimated rather than
ground-truth poses, and controlling the problem size, since every retained
point contributes a residual across all the keyframe pairs of the graph.
Robust estimators on the per-point residuals, a truncation band consistent
with the expected initial misalignment, and a keyframe-pair selection policy
derived from the loop closure machinery already available in RicoSLAM are the
natural starting points for this development;
\item \textbf{Broader baseline comparison and computational
characterization of RicoSLAM.}
The system-level evaluation of \chapref{chap:slam} compares RicoSLAM against
GMapping~\cite{grisetti:tro:2007} and
SLAM~Toolbox~\cite{macenski-jambrecic:joss:2021}, which supports the
conclusion that RicoSLAM outperforms these two baselines on the evaluated
sequences rather than a claim of general superiority over contemporary
2D~\gls{slam} systems. Extending the benchmark to additional systems, and to
further sequences from the literature, is required to generalize the
conclusions. A first step in this direction is already reported in the
manuscript in preparation~\cite{sousa:ral:2026} (\appxref{appx:ricoslam}),
which adds the Cartographer~\cite{hess:icra:2016} baseline, evaluated with
both its default and its real-time correlative scan matcher, together with
the total execution time of every framework on the same hardware.
Complementing that evaluation, the explicit measurement of the front-end and
back-end computation times, of the distance map precomputation cost against
the kd-tree~\cite{bentley:cacm:1975} search of classical \gls{icp}, and of
the memory footprint of each representation level remains to be performed.
\end{itemize}

\subsection{Involvement in the scientific community}\label{sec:conclusions:community}

In parallel with the research lines of this PhD~Thesis,
the candidate has been actively involved in the scientific
community through international research collaborations, the participation in
national funded scientific projects, and the co-supervision of BSc students at
\gls{feup}, \gls{isep}, and \gls{inesctec}. These activities provided
complementary perspectives to the research developed in this
PhD~Thesis and contributed to the mobile robotics team work on industrial
applications within \gls{criis}/\gls{iilab} at \gls{inesctec}.
This subsection consolidates these activities, organized into collaborations,
scientific projects, and supervision.

\subsubsection{Collaborations}

The international collaboration established with Prof.~Dr.~Giorgio~Grisetti from
Sapienza Università di Roma, materialized through a six-month research
visit between February and July 2024, shaped the
core contributions of the second part of the PhD~Thesis:
the distance map-based 2D~point cloud registration
framework introduced in \chapref{chap:distance-maps}, the RicoSLAM
framework presented in \chapref{chap:slam}, and the corresponding
distance map-based \gls{gba} formulations all originated from
discussions and joint development with Prof.~Dr.~Giorgio~Grisetti.
The visit also gave direct access to the srrg2\_solver~\cite{grisetti:robotics:2020}
optimization framework adopted as the back-end of RicoSLAM, and
provided the \emph{Cappero} and \emph{Kuka} sequences used to evaluate
the offline dynamics removal mechanism in
\secref{sec:slam:results:dynamics} on real-world data beyond the
IILABS~3D~\cite{ribeiro:access:2025} dataset and the
\gls{inesctec} industrial deployments.

At the national level, the candidate maintained an active collaboration
with the 5dpo robotics team~\cite{5dpo} at \gls{feup}, whose
three-wheeled omnidirectional platform was used as the experimental
platform for the corner-like feature detector reported in
\chapref{chap:features}. The collaboration with the 5dpo team
predates and continued throughout the PhD~Thesis, covering wheel
odometry calibration, pose control of omnidirectional robots, and
the development of a robotic framework in the scope of participating in multiple
editions of the Robot@Factory~4.0 competition at the RoboCup Portugal Open:
\begin{description}[leftmargin=1.25cm,style=nextline,font=\normalfont]
\item[\cite{sousa:icarsc:2021}] \fullcite{sousa:icarsc:2021}
\item[\cite{sousa:jint:2022}]   \fullcite{sousa:jint:2022}
\item[\cite{sousa:icarsc:2024}] \fullcite{sousa:icarsc:2024}
\item[\cite{lopes:icarsc:2025}] \fullcite{lopes:icarsc:2025}
\end{description}

\subsubsection{Scientific projects}

The PhD~Thesis was developed in articulation with a set of
applied research projects hosted at \gls{criis}/\gls{iilab} from \gls{inesctec}.
The primary alignment of this PhD~Thesis is with the
\textit{GreenAuto Agenda: Green Innovation for the Automotive Industry}~\cite{greenauto},
the main funding source of this PhD~Thesis through \gls{pps}18 -- \textit{3D~navigation system
for mobile robotic equipment} of \gls{wp}10 -- \textit{automated
logistics for the automotive industry}, as already discussed in
\secref{sec:intro:scope}. The Flowbotic real-world sequences in industrial-scale
environments considered in the dynamics removal experiment of
\secref{sec:slam:results:dynamics} are a consequence of this
articulation. In addition to GreenAuto, the candidate contributed to the
following scientific projects:
\begin{itemize}
\item \textbf{PRODUTECH~R3:} a national project on the
      reconfiguration, recovery, and resilience of the Portuguese
      manufacturing sector, within which the candidate contributed to
      the integration of mobile manipulators with RicoSLAM for tracking the
      robot pose during container unloading applications, 5G~communications
      for teleoperation of industrial \glspl{amr}, and pallet detection in
      intralogistics environments.
      This contribution led to the following publications:
  \begin{description}[leftmargin=1.25cm,style=nextline,font=\normalfont]
  \item[\cite{caldana:robot:2024}] \fullcite{caldana:robot:2024}
  \item[\cite{levin:robot:2024}] \fullcite{levin:robot:2024}
  \item[\cite{pacheco:icarsc:2025}] \fullcite{pacheco:icarsc:2025}
  \item[\cite{lopes:icara:2026}] \fullcite{lopes:icara:2026}
  \end{description}
\item \textbf{VINCI~7D:} a project focused on industrial printing-on-surfaces
      applications, in which the candidate contributed to a
      Gerber-format parser for printed-circuit-board layouts supporting bitmap
      and parametric representations, materialized in the publication below:
  \begin{description}[leftmargin=1.25cm,style=nextline,font=\normalfont]
  \item[\cite{sousa:access:2022}] \fullcite{sousa:access:2022}
  \end{description}
\end{itemize}

\subsubsection{Supervision}

Throughout the PhD~Thesis, the candidate co-supervised a number of
BSc students from \gls{feup} and \gls{isep} at \gls{inesctec}, on topics
related to the research lines of this work. Several of these supervisions
focused on modifying mobile platforms, sensorization, and \gls{slam} algorithms.
During the 2nd semester of the 2022/23 academic year (February to
June 2023), the candidate co-supervised BSc students from \gls{feup}
on the following topics, related to the PhD~Thesis's
focus on 2D laser-based localization and mapping for mobile robots:
\begin{itemize}
\item Ana Luísa Mota -- \emph{Localization and Mapping with the Hangfa
      Discovery~Q2 Mobile Platform}, exploring 3D~\gls{slam} with RGB-D sensors,
      while adopting the same platform modified in the scope of this
      PhD~Thesis~\cite{sousa:icarsc:2025} and used during the acquisition of the
      IILABS~3D~\cite{ribeiro:access:2025} dataset;
\item André Luís Pires -- \textit{Localization and Mapping with a
      Three-wheeled Omnidirectional Mobile Platform},
      addressing the same problem as Ana~Mota, but adopting the same platform
      of the 5dpo robotics team~\cite{5dpo} used as the experimental platform
      for the experimental evaluation in \chapref{chap:features};
\item José Miguel Guedes -- \textit{Trajectory Control of Omnidirectional Robots
      with the Hangfa Discovery~Q2 Mobile Platform}, focused on trajectory
      tracking control with the Hangfa~Discovery~Q2 platform,
      subsequently revised within this PhD~Thesis to acquire the
      IILABS~3D~\cite{ribeiro:access:2025} dataset;
\item Rodrigo de Vasconcelos e Miguel -- \textit{Motion Distortion Compensation
      on 2D~Laser Scanners Data}, a problem closely related to
      scan-distortion considerations that may affect \gls{slam} results.
\end{itemize}

In the same semester, the candidate also co-supervised a student from \gls{isep}:
\begin{itemize}
\item Paulo Ricardo Barbosa -- \textit{Test of 3D LiDAR SLAM Algorithms
      with the RoboSense RS-Helios-5515}, focused on the experimental
      validation of open-source 3D~\gls{slam} algorithms with the
      sensor referenced in the title, with this work predating the
      benchmark study by Ribeiro~\etal~\cite{ribeiro:access:2025}, a
      complementary contribution of this PhD~Thesis.
\end{itemize}

Finally, the candidate co-supervised \gls{bii} fellows at \gls{inesctec}
on topics that contributed to the project portfolio enumerated above
and, in some cases, materialized in the peer-reviewed publications
listed under the PRODUTECH~R3 project:
\begin{itemize}
\item Daniele da Mota Caldana --
      \textit{Pallet Detection in Intralogistics Environments},
      addressing pallet detection for autonomous handling within
      PRODUTECH~R3, with results contributing to the publication
      \cite{caldana:robot:2024};
\item Duarte Filipe Monteiro and David José Castro --
      \textit{Metallic Boxes Detection for the Automotive Sector},
      contributing to the metallic-box perception pipeline of the
      GreenAuto project that hosted this PhD~Thesis;
\item João Miguel Oliveira and Thiago Baldassarri Levin --
      \textit{Teleoperation of \glspl{amr} in 5G Networks},
      developing joystick-based teleoperation interfaces over
      5G~connectivity within PRODUTECH~R3, with results contributing
      to the publication \cite{levin:robot:2024}.
\end{itemize}
